% Multimodal Speech Emotion Recognition with RoBERTa and wav2vec2:
% A Two-Dataset Study of Modality Complementarity
%
% Author: Shahnawaz Khan
% Target venue: 4-page short paper, Interspeech / ICASSP workshop style.
% Compile with: pdflatex paper.tex (after copying into a proper LaTeX template
%               of your choice — e.g. IEEEtran, interspeech2025.sty, etc.)
%
% This file is written in a venue-agnostic single-column draft format.
% Section headings, citations, and figure references are placeholder-ready.

\documentclass[a4paper,10pt]{article}

\usepackage[utf8]{inputenc}
\usepackage{microtype}
\usepackage{booktabs}
\usepackage{hyperref}
\usepackage{graphicx}
\usepackage{url}
\usepackage{xcolor}

\title{Multimodal Speech Emotion Recognition with RoBERTa and wav2vec2:\\
       A Two-Dataset Study of Modality Complementarity}
\author{Shahnawaz Khan\\
        Independent Researcher\\
        \texttt{contact@skakarh.com}\\
        \texttt{github.com/ShahnawazKakarh}}
\date{June 2026}

\begin{document}
\maketitle

\begin{abstract}
We train and compare text-only (RoBERTa), audio-only (wav2vec2 / WavLM), and
multimodal cross-attention models for Speech Emotion Recognition (SER) on
two standard datasets: RAVDESS (acted lab speech, 8 emotions, 1\,440 clips)
and MELD (\emph{Friends} TV dialogue, 7 emotions, 13\,708 utterances). Using
a single shared architecture and identical training pipeline, multimodal
fusion improves weighted F1 by \textbf{+6.9 percentage points} over
audio-only on RAVDESS (speaker-independent split) but \emph{underperforms}
text-only by \textbf{1.9 percentage points} on MELD. We attribute the
divergence to modality complementarity: clean acted speech with degenerate
textual content (RAVDESS has only two carrier sentences) rewards fusion,
while noisy multi-party conversational audio paired with rich semantic text
(MELD) does not. The framework, configurations, trained checkpoints, and a
live HuggingFace Spaces demo are publicly released to support reproducibility.
\end{abstract}

\section{Introduction}

Speech Emotion Recognition is a long-standing problem with strong unimodal
baselines from both the speech and NLP communities, and a recurring claim
in recent multimodal work that text-plus-audio fusion is uniformly
beneficial \cite{poria2017,zadeh2018,hazarika2020}. We test that claim
under controlled conditions: same architecture, same training pipeline, two
datasets chosen for opposite modality profiles. The results challenge the
``multimodal is always better'' framing and reproduce a finding more
familiar from CMU-MOSEI literature \cite{zadeh2018,hazarika2020} on a
purely SER setting.

Our contributions are:
\begin{enumerate}
\item A reproducible PyTorch Lightning codebase\footnote{%
\url{https://github.com/ShahnawazKakarh/speech-emotion-recognition-transfer-learning}}
implementing text-only, audio-only, and cross-attention multimodal SER
with a shared training and evaluation pipeline.
\item Speaker-independent baselines on RAVDESS with per-class F1
breakdown, quantifying the ``speaker-leakage premium'' inherent in random
splits ($\sim$13~pp WF1 inflation).
\item MELD baselines (text / audio / multimodal) that match or exceed the
published literature range, with honest reporting of the audio-only
class-collapse failure mode (UF1 = 0.153).
\item A live, public demo running a pretrained baseline\footnote{%
\url{https://huggingface.co/spaces/Shahnawazkakarh/speech-emotion-recognition}}.
\end{enumerate}

\section{Method}

\subsection{Architecture}

All three models share a common Lightning module that exposes text and
audio encoders, an optional fusion module, and a classifier head.

\textbf{Text encoder.} \texttt{roberta-base} with mean-pooling over the
\texttt{[CLS]} token's contextual representation. For MELD, we prepend up
to two prior utterances from the same dialogue, separated by \texttt{[SEP]},
to give the encoder conversational context.

\textbf{Audio encoder.} \texttt{facebook/wav2vec2-base} (RAVDESS) or
\texttt{microsoft/wavlm-base} (MELD), with frozen CNN feature extractor and
selectively-frozen transformer layers. Output is the mean-pooled final
hidden state.

\textbf{Fusion module.} Two-layer cross-attention block where text tokens
attend over audio frames (and vice versa), followed by a feed-forward
projection to a 512-dimensional fused representation. We use 8 attention
heads and dropout 0.2.

\textbf{Classifier.} A two-layer MLP head over the (pooled) representation
with dropout 0.2 and softmax output.

\subsection{Training}

Cross-entropy with inverse-frequency class weights. AdamW with linear
warmup-and-decay schedule. Learning rate 2e-5 for unimodal runs (after
verifying that the 1e-4 pretraining rate causes representation collapse;
see \S\ref{sec:hparams}) and 5e-5 for multimodal. Maximum 10 epochs with
EarlyStopping (patience 3) on validation WF1. All checkpoints select by
best validation WF1.

Apple Silicon MPS backend throughout. For the 221\,M-parameter multimodal
MELD configuration, we use \texttt{batch\_size=2} with
\texttt{accumulate\_grad\_batches=8} (effective batch 16) and audio clips
truncated to 5 seconds to fit within the unified-memory budget.

\subsection{Datasets and splits}

\textbf{RAVDESS.} 1\,440 clips, 8 emotion classes, 24 actors, 2 carrier
sentences. We use the published speaker-independent (SI) split: actors
1--18 train, 19--20 validation, 21--24 test (240 test samples). For
comparison we also report a random 70/10/20 split. All audio is 16\,kHz mono.

\textbf{MELD.} 13\,708 utterances from \emph{Friends}, 7 emotions, with
official train / dev / test partitions (9\,989 / 1\,109 / 2\,610). We
extract 16\,kHz mono audio from the original \texttt{.mp4} videos using
ffmpeg. Three utterances (one per split) had corrupted source videos
producing zero-byte \texttt{.wav}; these are filtered at loader
initialization with a logged warning. The resulting test set is 2\,609
utterances.

\section{Results}

\subsection{Main results}

Table~\ref{tab:main} reports test-set weighted F1 (WF1), unweighted F1
(UF1, equivalent to UAR), and accuracy for all three architectures on both
datasets.

\begin{table}[t]
\centering
\caption{Test-set performance on RAVDESS (speaker-independent split) and
MELD (official split). Best per dataset in \textbf{bold}.}
\label{tab:main}
\small
\begin{tabular}{lccc}
\toprule
\textbf{Model} & \textbf{WF1} & \textbf{UF1} & \textbf{Acc} \\
\midrule
\multicolumn{4}{l}{\emph{RAVDESS (SI split)}} \\
Text-only RoBERTa & 0.031 & 0.029 & 0.133 \\
Audio-only wav2vec2 & 0.659 & 0.631 & 0.667 \\
\textbf{Multimodal} (RoBERTa + wav2vec2) & \textbf{0.728} & \textbf{0.731} & \textbf{0.729} \\
\midrule
\multicolumn{4}{l}{\emph{MELD (official split)}} \\
Audio-only WavLM & 0.357 & 0.153 & 0.416 \\
Multimodal (RoBERTa + WavLM) & 0.590 & 0.404 & 0.597 \\
\textbf{Text-only RoBERTa (+ ctx)} & \textbf{0.609} & \textbf{0.459} & \textbf{0.593} \\
\bottomrule
\end{tabular}
\end{table}

The headline observation is the divergence between datasets: multimodal
fusion improves WF1 by \textbf{+6.9~pp} over the best unimodal baseline on
RAVDESS, but \emph{decreases} WF1 by \textbf{1.9~pp} on MELD. Across both
datasets, the same architecture and pipeline produce opposite verdicts.

\subsection{Per-class behaviour}

Table~\ref{tab:perclass} shows per-class F1 for the multimodal model
against its strongest unimodal competitor on each dataset.

\begin{table}[t]
\centering
\caption{Per-class F1: multimodal vs.\ strongest unimodal baseline on each
dataset. Positive $\Delta$ = multimodal improves; negative $\Delta$ =
multimodal hurts.}
\label{tab:perclass}
\small
\begin{tabular}{lccr}
\toprule
\textbf{Class} & \textbf{Multimodal} & \textbf{Unimodal} & \textbf{$\Delta$} \\
\midrule
\multicolumn{4}{l}{\emph{RAVDESS SI (Multimodal vs.\ Audio-only)}} \\
neutral   & \textbf{0.778} & 0.211 & \textbf{+56.7 pp} \\
calm      & 0.776 & 0.691 & +8.5 pp \\
happy     & 0.603 & 0.590 & +1.3 pp \\
sad       & 0.607 & 0.508 & +9.9 pp \\
angry     & 0.794 & 0.828 & --3.4 pp \\
fearful   & 0.745 & 0.720 & +2.5 pp \\
disgust   & 0.750 & 0.787 & --3.7 pp \\
surprised & 0.795 & 0.714 & +8.0 pp \\
\midrule
\multicolumn{4}{l}{\emph{MELD (Multimodal vs.\ Text-only)}} \\
neutral   & 0.755 & 0.733 & +2.1 pp \\
joy       & 0.582 & 0.587 & --0.5 pp \\
sadness   & 0.293 & 0.353 & --5.9 pp \\
anger     & 0.372 & 0.487 & \textbf{--11.5 pp} \\
surprise  & 0.537 & 0.576 & --3.9 pp \\
fear      & 0.059 & 0.205 & \textbf{--14.7 pp} \\
disgust   & 0.229 & 0.273 & --4.3 pp \\
\bottomrule
\end{tabular}
\end{table}

On \textbf{RAVDESS}, multimodal helps every class except angry and disgust,
where audio alone already achieves high F1 ($\geq$0.79). The neutral class
sees a dramatic +56.7~pp gain: audio-only fails on neutral utterances from
unseen speakers (F1 = 0.21, recall = 0.125), and the text branch -- which
predicts ``calm'' constantly -- supplies enough disambiguation prior to
recover them.

On \textbf{MELD}, multimodal helps only the dominant neutral class (+2.1~pp)
and hurts every other class, with the largest drops on fear (--14.7~pp),
anger (--11.5~pp), and sadness (--5.9~pp). These are precisely the emotions
that depend most on lexical / semantic content, which the text branch
captures well alone; fusing with weak audio degrades the representation.

The audio-only MELD model exhibits \emph{class collapse}: F1 = 0.000 for
surprise, fear, and disgust -- three of seven classes contributing nothing
to predictions. Best validation WF1 occurred at epoch 0, with all
subsequent training raising training accuracy but not validation
performance. This pattern is consistent with published WavLM-based MELD
baselines and reflects the difficulty of acoustic-only discrimination on
noisy, severely-imbalanced TV-show audio.

\subsection{Hyperparameter sensitivity}
\label{sec:hparams}

A preliminary RAVDESS audio-only run with learning rate 1e-4 (the
wav2vec2 pre-training rate) diverged: model unlearned representations
within 4 epochs, EarlyStopping triggered at WF1 0.27. Reducing to 2e-5
with the bottom 8 of 12 transformer layers frozen reached WF1 0.796 on
random split / 0.659 on SI -- a 3$\times$ performance gap from a 5$\times$
LR difference. We adopted 2e-5 across all subsequent audio fine-tuning.

\subsection{Speaker-leakage premium}

For sanity, we additionally ran RAVDESS under a random 70/10/20 split
(same actors in train and test). Multimodal random WF1 was 0.858 vs.\
0.728 SI -- a \textbf{13~pp inflation}. We flag this throughout as the
\emph{speaker-leakage premium}: a substantial fraction of published
random-split RAVDESS WF1 reports describe a different (easier) problem
than speaker-independent evaluation.

\section{Discussion}

\subsection{When does multimodal help?}

The two-dataset comparison supports a simple rule: \emph{multimodal fusion
helps when both modalities carry independent, non-degenerate signal; it
hurts when one modality is strong and the other is noise.}

RAVDESS audio is rich (acted prosody); RAVDESS text is degenerate (two
fixed sentences). Multimodal wins. MELD text is rich (real dialogue
content); MELD audio is degraded (laugh tracks, multi-speaker scenes,
extreme class imbalance). Multimodal loses.

This aligns with prior multimodal fusion findings on CMU-MOSEI sentiment
analysis \cite{zadeh2018,hazarika2020}, where text typically dominates and
audio adds limited value, but rarely with controlled within-codebase
demonstration on emotion data specifically.

\subsection{Practical implications}

For practitioners deploying SER systems, the implication is to evaluate
unimodal baselines on the target domain before committing to multimodal
inference cost. A 1.9~pp text-only advantage at $\sim$45\% the parameter
count (124\,M vs 221\,M) and $\sim$30\% the inference latency is a
meaningful efficiency win that the multimodal framing tends to overlook.

\subsection{Limitations}

(1) Both datasets are English-only. We have not tested cross-lingual
transfer, which is a planned follow-up using XLM-R and multilingual
wav2vec2 on a self-recorded Urdu/Hindi/Punjabi corpus.

(2) IEMOCAP, the third standard SER dataset, requires a USC SAIL license
we have not yet received. Results on IEMOCAP would be needed to claim the
finding generalizes beyond RAVDESS / MELD.

(3) The cross-attention fusion is one of many possible architectures.
Tensor fusion, gated multimodal units \cite{arevalo2017}, and modality-
invariant/specific representations (MISA, \cite{hazarika2020}) might
exhibit different sensitivity to noisy modalities.

(4) MELD audio-only class-collapse may be partially attributable to
fine-tuning hyperparameters or to the WavLM-base capacity ceiling; we did
not run an exhaustive sweep.

\section{Conclusion}

A controlled, single-codebase comparison of text-only, audio-only, and
multimodal cross-attention SER reveals that the answer to ``does
multimodal fusion help?'' is dataset-dependent. Multimodal wins on
RAVDESS by +6.9~pp WF1; it loses on MELD by 1.9~pp WF1. The determining
factor is modality complementarity: fusion is beneficial only when each
modality contributes independent, non-noisy signal. Code, configs, and a
live demo are released for inspection and extension.

\section*{Reproducibility}

All experiments are reproducible from the released repository
(\url{https://github.com/ShahnawazKakarh/speech-emotion-recognition-transfer-learning}).
The repository includes:

\begin{itemize}
\item YAML configs for all 9 experiments (3 RAVDESS-random, 3 RAVDESS-SI,
3 MELD), with seeded splits.
\item Training scripts (\texttt{src/train.py}) and evaluation scripts
(\texttt{src/evaluate.py}) using PyTorch Lightning.
\item Data preparation scripts (\texttt{scripts/prepare\_meld.sh},
\texttt{scripts/download\_ravdess.sh}) with integrity checks and resume
support.
\item A live HuggingFace Space\footnote{%
\url{https://huggingface.co/spaces/Shahnawazkakarh/speech-emotion-recognition}}
running a pretrained baseline for instant inference.
\item Full per-class metrics and confusion matrices in
\texttt{results/results.md}.
\end{itemize}

\section*{Future work}

The next research arc, currently in planning, extends this framework to
\textbf{cross-lingual SER}: applying English-trained models to lower-resource
languages (Urdu, Hindi, Punjabi) via XLM-R for text and multilingual
wav2vec2 for audio, with a small self-recorded target-language evaluation
corpus. Target venues include the Interspeech / ICASSP workshop track and
IEEE Transactions on Audio, Speech, and Language Processing.

\begin{thebibliography}{9}

\bibitem{poria2017} S.\ Poria, E.\ Cambria, R.\ Bajpai, A.\ Hussain,
``A review of affective computing: From unimodal analysis to multimodal
fusion,'' \emph{Information Fusion}, vol.\ 37, pp.\ 98--125, 2017.

\bibitem{zadeh2018} A.\ Zadeh et al., ``Multimodal language analysis in
the wild: CMU-MOSEI dataset and interpretable dynamic fusion graph,'' in
\emph{Proc.\ ACL}, 2018.

\bibitem{hazarika2020} D.\ Hazarika, R.\ Zimmermann, S.\ Poria, ``MISA:
Modality-invariant and -specific representations for multimodal sentiment
analysis,'' in \emph{Proc.\ ACM MM}, 2020.

\bibitem{arevalo2017} J.\ Arevalo et al., ``Gated multimodal units for
information fusion,'' in \emph{ICLR Workshop}, 2017.

\bibitem{livingstone2018} S.\ R.\ Livingstone, F.\ A.\ Russo, ``The Ryerson
Audio-Visual Database of Emotional Speech and Song (RAVDESS),''
\emph{PLOS ONE}, 13(5), 2018.

\bibitem{poria2019} S.\ Poria et al., ``MELD: A Multimodal Multi-Party
Dataset for Emotion Recognition in Conversations,'' in \emph{Proc.\ ACL},
2019.

\bibitem{baevski2020} A.\ Baevski et al., ``wav2vec 2.0: A framework for
self-supervised learning of speech representations,'' in \emph{Proc.\
NeurIPS}, 2020.

\bibitem{chen2022} S.\ Chen et al., ``WavLM: Large-scale self-supervised
pre-training for full stack speech processing,'' \emph{IEEE Journal of
Selected Topics in Signal Processing}, 2022.

\bibitem{liu2019} Y.\ Liu et al., ``RoBERTa: A robustly optimized BERT
pretraining approach,'' \emph{arXiv:1907.11692}, 2019.

\end{thebibliography}

\end{document}
